\let\cleardoublepage\clearpage
\part{Производственные и обрабатывающие отрасли}
\chapter{Пищевая технология}
\ID{IRSTI 65.63.91}{}

\begin{header}
\swa{T.Ch.~Tultabaeva, G.N.~Zhakupova, К.К.~Makangali, А.Т.~Sagandyk, А.K.~Muldasheva, К.Sh.~Dairova}{A RESEARCH OF THE COMPOSITION AND FUNCTIONAL PROPERTIES OF BUTTERMILK CONCENTRATE OBTAINED THROUGH ULTRAFILTRATION}

T.Ch.~Tultabaeva,
G.N.~Zhakupova,
К.К.~Makangali,
А.Т.~Sagandyk,
А.K.~Muldasheva,
К.Sh.~Dairova\envelope
\end{header}

\begin{affil}
«S.Seifullin Kazakh Agro Technical Research University», Astana, Kazakhstan

\corrauthor{Corresponding author: dairova111@mail.ru}
\end{affil}

Under the growing demand for functional and diet-oriented dairy
products, this study is devoted to the development of a fermented
beverage technology enriched with high-quality protein through the
rational utilization of secondary dairy resources -- specifically,
buttermilk subjected to ultra\-filtration. The objective was to
investigate the composition, physicochemical parameters, and sensory
attributes of buttermilk concentrates obtained under varying
ultrafiltration conditions, and to assess their technological viability
within multicomponent dairy matrices.

Experimental trials revealed the impact of processing temperature
(20--60\,°C) on changes in pH, titra\-table acidity, fat and protein
contents, as well as total solids, in both retentate and permeate
fractions. Ultra\-filtration proved highly effective at retaining protein
and fat components with minimal loss, producing a concentrate of notable
nutritional and biological value.

The resulting buttermilk concentrate was incorporated into dairy blends
containing various ratios of cow's and goat's milk. Across four
formulated variants, a consistent trend was observed: increasing the
proportion of buttermilk concentrate led to a corresponding rise in
total solids (from 18.25\% to 35.50\%) and protein content (from 6.01\%
to 14.00\%).

Thus, buttermilk concentrate derived via ultrafiltration emerges as a
promising ingredient for the next generation of functional dairy
products - capable of enhancing protein content, contributing to
desirable texture formation, and reducing overall fat levels in the
final formulation.

{\bfseries Key words:} buttermilk, functional properties, permeate,
retentate, ultrafiltration.

\begin{header}
УЛЬТРАСҮЗГІЛЕУ ӘДІСІМЕН АЛЫНҒАН ІРКІТ КОНЦЕНТРАТЫНЫҢ ҚҰРАМЫ МЕН ҚАСИЕТТЕРІН ЗЕРТТЕУ

Т.Ч.~Тултабаева,
Г.Н.~Жакупова,
Қ.Қ.~Мақанғали,
А.Т.~Сағандық,
А.Х.~Мулдашева,
Қ.Ш.~Даирова\envelope
\end{header}

\begin{affil}
«С.Сейфуллин атындағы Қазақ агротехникалық зерттеу университеті», Астана, Қазақстан,

e-mail: dairova111@mail.ru
\end{affil}

Функционалдық және диеталық сүт өнімдеріне сұраныстың артуы жағдайында
бұл зерттеу жоғары сапалы ақуызға бай ферменттелген сусындар
технологиясын әзірлеуге бағытталған. Бұл мақсатқа сүт өнеркәсібінің
екінші реттік ресурстарын, атап айтқанда, ультрасүзгіден өткен іркітті
тиімді пайдалану арқылы қол жеткізілді. Зерттеудің мақсаты -- іркіт
концентраттарының құрамы, физика-химиялық және органолептикалық
қасиеттерін әртүрлі ультрасүзу режимдерінде зерттеу, сондай-ақ оларды
көп компонентті сүт қоспалары құрамында технологиялық тұрғыдан бағалау
болды.

Жүргізілген тәжірибелер температураның (20--60\,°C) рН, титрленетін
қышқылдық, май, ақуыз және құрғақ заттардың мөлшеріне әсерін анықтады.
Ультрасүзу процесі ақуыздар мен майларды ең аз жоғалту арқылы тиімді
түрде ұстап, жоғары тағамдық және биологиялық құндылығы бар концентрат
қалыптастыратыны дәлелденді.

Іркіт концентраты әртүрлі мөлшерде сиыр және ешкі сүтімен араласқан сүт
негізіне енгізілді. Төрт тәжірибелік қоспада іркіт концентратының үлесі
артқан сайын құрғақ заттар мөлшері (18,25\%-дан 35,50\%-ға дейін) және
ақуыз мөлшері (6,01\%-дан 14,00\%-ға дейін) тұрақты түрде өсті.

Осылайша, ультрасүзу әдісімен алынған іркіт концентраты функционалдық
сүт өнімдерін жасауға арналған келешегі зор ингредиент ретінде
қарастырылады. Ол өнімнің ақуыздық құрамын байытып, құрылым түзілуіне
ықпал етіп, дайын өнімнің майлылығын төмендетуге мүмкіндік береді.

{\bfseries Түйін сөздер:} іркіт, функционалдық қасиеттері, пермеат,
ретентат, ультрасүзгілеу.

\begin{header}
ИССЛЕДОВАНИЕ СОСТАВА И СВОЙСТВ КОНЦЕНТРАТА ПАХТЫ, ПОЛУЧЕННОГО МЕТОДОМ УЛЬТРАФИЛЬТРАЦИИ

Т.Ч.~Тултабаева,
Г.Н. Жакупова,
К.К. Маканғали,
А.Т.~Сагандык,
А.Х.~Мулдашева,
К.Ш.~Даирова\envelope
\end{header}

\begin{affil}
«Казахский агротехнический исследовательский университет им. С.Сейфуллина», Астана, Казахстан,

e-mail: dairova111@mail.ru
\end{affil}

В условиях растущего спроса на функциональные и диетические молочные
продукты данное исследование направлено на разработку технологии
ферментированных напитков с повышенным содержанием полноценного белка за
счёт рационального использования вторичных ресурсов молочной
промышленности -- в частности, пахты, подвергнутой ультрафильтрации.
Целью работы стало изучение состава, физико-химических и
органолептических свойств концентратов пахты, полученных при различных
режимах ультрафильтрационной обработки, а также оценка их
технологического потенциала в составе многокомпонентных молочных смесей.

Экспериментально установлено влияние температуры (20--60\,°С) на
изменение pH, титруемой кислотности, содержания жира, белка и сухих
веществ в ретентате и пермеате. Установлено, что ультрафильтрация
обеспечивает эффективное удержание белковых и жировых компонентов при
минимальных потерях, формируя концентрат с высокой пищевой и
биологической ценностью.

Концентрат пахты использован в составе молочной основы с различным
соотношением коровьего и козьего молока. В четырёх опытных вариантах
смеси наблюдалось закономерное увеличение массовой доли сухих веществ
(от 18,25\% до 35,50\%) и белка (от 6,01\% до 14,00\%) при возрастании
доли концентрата пахты.

Таким образом, концентрат пахты, полученный методом ультрафильтрации,
является перспективным ингредиентом для функциональных молочных
продуктов, обеспечивающим белковое обогащение, текстурообразование и
снижение жирности готового продукта.

{\bfseries Ключевые слова:} пахта, функциональные свойства, пермеат,
ретентат, ультрафильтрация.

\begin{multicols}{2}
{\bfseries Introduction.} One of the strategic priorities in the current
development of the dairy industry is the formulation and implementation
of technologies for functional food products that support rational
nutrition, sustain physical and cognitive performance, and enhance the
human immune system. In this context, a particularly relevant objective
is the production of fermented dairy products enriched with high-quality
proteins.

Although the market already offers a broad range of fermented dairy
items, there remains a notable shortage of products with reduced fat
content, those derived from secondary dairy raw materials --
particularly buttermilk - and specialized dairy products tailored for
individuals with health conditions such as diabetes or lactose
intolerance.

At present, new categories of fermented dairy products are being
developed with pronounced therapeutic and prophylactic properties. The
functionalization of such products can be achieved by utilizing milk
from various agricultural species, applying membrane processing
technologies, and adopting low-waste or zero-waste production systems in
the development of fermented dairy products.

The dairy industry faces a persistent shortage of primary raw materials,
which brings to the forefront the issue of rational utilization of
secondary dairy resources such as skim milk, buttermilk, and whey.

Comprehensive utilization of all milk components can significantly
diversify the product range and positively impact the technical and
economic performance indicators, as well as labor productivity in dairy
enterprises. Given their considerable volume and high nutritional and
biological value, the collection and processing of skim milk,
buttermilk, and whey become not only justified but essential {[}1{]}.

In countries with advanced dairy sectors, buttermilk is widely employed
in the manufacture of various dietary natural and cultured beverages,
including numerous variants with flavorings and additives, as well as an
ingredient in food formulations. To extend shelf life and more
efficiently utilize available buttermilk resources - especially during
peak milk processing periods for butter production-\/-- the manufacture
of buttermilk concentrates is considered preferable.

{\bfseries Materials and methods.} In the course of the study, various
dairy substrates were utilized as primary samples, chosen for their
distinct protein-fat compositions and functional properties. The main
object of analysis was buttermilk taken from the ЕС "Astana-Onim" - a
by-product of butter production characterized by reduced fat content and
enriched with polar lipids as well as biologically active peptides.

By means of ultrafiltration of buttermilk using semipermeable membranes,
two fractions were obtained:

- permeate, which predominantly contains low-molecular-weight compounds
such as lactose, mineral salts, and soluble vitamins,

- retentate, enriched with proteins, phospholipids, and
high-molecular-weight components {[}2{]}.

The schematic representation of the ultrafiltration setup is shown in
Figure 1.
\end{multicols}

\begin{figs}[Fig.1 - The process of ultrafiltration of buttermilk\\\normalfont{{\em where, 1 - Sample / diafiltration reservoir, 2 - Filtrate versel, 3 -
Pressure indicator, 4 - Pink flow restrictor, 5 - Masterflex pump, 6 -
Masterflex easy load pump head.}}]
    \fig[0.45\textwidth]{p/image7}[The scheme of the ultrafiltration plant]
    \fig[0.45\textwidth]{p/image8}[The process
of ultrafiltration of buttermilk]
\end{figs}

\begin{multicols}{2}
In the VivaFlow ultrafiltration system, the process works as follows:
the feed solution is placed in the diafiltration reservoir (1), from
where it is pumped to the membrane using the Masterflex pump (5) with
the Easy Load pump head (6). The membrane separates the liquid into
concentrate and filtrate, with the filtrate collected in a separate
vessel (2). System pressure is monitored by the pressure indicator (3),
and the flow is controlled by the pink flow restrictor (4), ensuring
stable operation and efficient separation of the solution components.

Along with buttermilk and its derivatives, Simmental cow's milk from LLP
«Moya ferma» and Saanen goat's milk from LLP «Arka Zerenda» were
utilized in the study to enable a comparative analysis of
species-specific differences in casein micelle structure and lipid
composition.

Furthermore, the investigation included a buttermilk concentrate,
produced by partial dehydration, aimed at modeling highly concentrated
dairy systems.

Finally, the research incorporated blends of various dairy components,
specifically prepared to explore synergistic effects between proteins,
fats, and polar lipids within multiphase dairy matrices.

Active acidity (pH) was determined in accordance with «GOST 26188-84 -
Food products. Method for pH determination», using an Atago pH meter,
which ensures reliable and precise assessment of hydrogen ion
concentration in the tested samples {[}3{]}.

Titratable acidity (°T) was assessed using the acid-base titration
method in accordance with «GOST 3624-92 - Milk and dairy products.
Method for titratable acidity determination», providing an integral
measure of all titratable acids present in the product {[}4{]}.

Mass fraction of fats (\%) was measured using Fourier-transform
near-infrared (FT-NIR) spectroscopy on a Tango-R spectrometer, a rapid
and non-destructive analytical method.

Mass fraction of proteins (\%) was also determined using FT-NIR
spectroscopy on the Tango-R device, which captures absorption spectra
characteristic of protein molecular structures, serving as an effective
alternative to classical chemical assays.

Mass fraction of dry substances (\%) was determined according to «GOST
3626-73 - Milk and dairy products. Method for determination of dry
matter content», by means of the refractometric method, which relates
the refractive index of the sample to its concentration of dissolved
solids {[}5{]}.

The sensory analysis was carried out in accordance with GOST R ISO
22935-2-2011 «Milk and milk products. Sensory analysis. Part 2.
Recommended methods for sensory evaluation» {[}6{]}. The tasting was
attended by teachers and students of KATRU, a total of 7 degustators.
The assessment was carried out in a tasting room at a temperature of (20
± 2) ° C and diffused daylight, eliminating the influence of foreign
odors and glare. The samples were presented in the same container,
randomly, in increasing order of concentration of dry substances, with
an interval of 2-3 minutes between assessments to restore taste buds.
The sensory evaluation was carried out using a five-point scale, where 5
points corresponded to samples with the most desirable organoleptic
characteristics and 1 point to those with pronounced defects or
deviations from the standard. The results of the sensory evaluations
were statistically processed by calculating the mean values and standard
deviations.

{\bfseries Results and discussion.} One of the promising avenues for the
valorization of whey and buttermilk is membrane filtration, a technology
that enables the selective separation of the product into two distinct
fractions: permeate and retentate. The present study investigates the
temperature-dependent changes (in the range of 20-60\,°C) in the
physicochemical properties of buttermilk and its ultrafiltration-derived
fractions {[}7{]}.

Ultrafiltration was performed using the Vivaflow 50/50R/200 system, and
the technical specifications of the equipment are presented in Table 1.

The raw material - buttermilk - was poured into a product tank, from
which it was transferred into a membrane module via a membrane pump.
Within the module, the feed was separated into retentate (concentrate)
and filtrate {[}8{]}. The ultrafiltration pressure was regulated using a
needle valve and monitored by a manometer, while temperature was
maintained by a heat exchanger. During the process, the retentate was
recirculated back into the product tank, and the filtrate was collected
in a separate container.

It is well established that increasing temperature enhances filtration
rate, whereas membrane selectivity remains largely unaffected. However,
elevating the temperature beyond 40\,°C can lead to the precipitation of
soluble calcium and magnesium salts, resulting in pore blockage and a
reduction in filtration efficiency {[}9{]}.
\end{multicols}

\tcap{Table 1 - Technical Specifications of the Ultrafiltration System Vivaflow 50/50R/200}
\begin{longtblr}[
  label = none,
  entry = none,
]{
  cells = {c},
  hlines,
  vlines,
}
Overall dimensions (length/height/width), mm & 107/84/25      \\
Channel dimensions (width/height), mm        & 15/0,3         \\
Effective membrane area, cm²                 & 50             \\
Minimum recirculation volume, mL             & 10             \\
Membrane material                            & ТРХ (PMF)      \\
Manufacturer                                 & SARTORIUS (UK) 
\end{longtblr}

\fig[0.75\textwidth]{p/image9}[Fig.2 - Comparative characteristics of buttermilk
processing methods]

\begin{multicols}{2}
Buttermilk concentrates obtained through membrane separation
technologies represent valuable ingredients for the food industry. Their
composition and functional properties vary depending on the filtration
method - whether ultrafiltration or nanofiltration {[}10{]}. Given the
absence of a specific GOST standard regulating the physicochemical
properties of ultrafiltered buttermilk concentrate, a comparative
evaluation was carried out based on the existing GOST 34354-2017 «Dairy
products. Buttermilk concentrate. Technical specifications» (which
pertains to condensed buttermilk), as well as the study by Bobrova,
which investigates the characteristics of nanofiltered buttermilk
{[}11{]}.

Figure 2 presents сomparative characteristics of buttermilk processing
methods.

According to GOST requirements for condensed buttermilk, titratable
acidity must not exceed 21 °T. In both membrane-processed samples this
level was surpassed: 32.0 ± 1.0 °T in the nanofiltration sample and 25.6
± 1.0 °T in the ultrafiltration sample. The increase likely reflects
active microbial development during concentration or insufficient
thermal stabilization of the raw material. Since temperature rises were
observed in both processes, membrane filtration may intensify acidity
either through microbiological activity or through concentration
effects. Notably, the higher acidity of the nanofiltration sample may
accelerate spoilage compared to ultrafiltration.

Fat content in both concentrates remained within the GOST-specified
range (0.5--1.5\%), amounting to 0.78 ± 0.08\% after nanofiltration and
0.72 ± 0.08\% after ultrafiltration, confirming effective lipid
retention. The slight decrease in the ultrafiltration sample may be
linked to the coarser membrane, which can allow finer lipid fractions to
pass.

Protein content differed markedly: the nanofiltration concentrate
reached 7.03 ± 0.30\%, exceeding the normative range (2.5-4.0\%), while
the ultrafiltered sample contained 4.00 ± 0.20\%, close to the upper
GOST limit. This distinction highlights the higher selectivity of
nanofiltration, which retains both high-molecular proteins and smaller
peptides.

According to GOST, the dry matter content of condensed buttermilk should
be 30-50\% (or above 90\% for dried products). In these
membrane-concentrated samples, whey removal did not occur; concentration
was achieved only by moisture reduction. As a result, dry matter reached
20.10 ± 0.50\% after nanofiltration and 17.7 ± 0.50\% after
ultrafiltration.

Overall, the results show that membrane concentration effectively
preserves proteins and fats. Ultrafiltration stands out as a practical
and economically advantageous method, providing substantial retention of
proteins and dry matter and making it suitable for the production of
functional food ingredients, while maintaining lower acidity compared to
nanofiltration.

The experiment was conducted with repeated measurements: each
temperature condition (20, 30, 40, 50, and 60 °C) was tested in three
parallel replicates, allowing for reliable mean values to be obtained
and for statistical differences between the conditions to be assessed.
Table 2 provides a comparative summary of key physicochemical parameters
of buttermilk concentrates produced via different membrane methods,
alongside the reference values from GOST 34354--2017 {[}12{]}.

In the control buttermilk sample, the active acidity was measured at pH
7.14, indicating a mildly alkaline medium. Upon thermal treatment, a
gradual decline in pH was observed, most notably in the permeate
fraction. At 60\,°C, the pH of the permeate decreased to 6.42,
suggesting intensified acid-forming reactions. In contrast, the
retentate exhibited a narrower pH fluctuation range (from 6.84 to 6.56),
likely attributable to its higher protein content, which provides
buffering capacity.

Titratable acidity increased with rising temperature, peaking at 40\,°C
with values of 30\,±\,2\,°T in the retentate and 27.0\,°T in the
permeate. A subsequent decline in acidity beyond 50\,°C may reflect the
thermal inactivation of microbial populations and modifications in the
protein--mineral matrix.

The mass fraction of fat in the permeate remained consistently low
(0.06--0.1\%), confirming the efficiency of membrane separation in
removing lipid components from the aqueous phase. Meanwhile, the
retentate demonstrated a gradual increase in fat content with
temperature elevation -- from 0.70 to 0.95\% - possibly due to fat
globule densification and enhanced retention within the concentrated
phase. The most pronounced changes were observed in the distribution of
protein between the ultrafiltration fractions. In the permeate, the mass
fraction of protein decreased from 1.90\% to 1.44\% at 40\,°C, followed
by a slight increase at higher temperatures. In contrast, the retentate
maintained a relatively stable protein content, fluctuating between
3.80\% and 4.00\%, indicating efficient retention of
high-molecular-weight compounds.

The mass fraction of dry matter in the permeate showed a steady decline
from 5.29\% to 4.94\%, whereas in the retentate, dry matter content
accumulated progressively, ranging from 10.29\% to a maximum of
15.38--17.7\%. These data underscore the effectiveness of thermal
concentration at 30\,°C, where the highest dry matter content (17.7\%)
was recorded, suggesting this temperature regime as the most favorable
for obtaining highly concentrated dairy products.

According to the experimental results presented in Table 2, the maximum
protein yield (4\%) was achieved at 30\,°C. The concentration process
was carried out at 30\,°C and a pressure of 2.5\,MPa. From an initial
volume of 500\,mL of buttermilk, the process yielded 400\,mL of whey and
100\,mL of buttermilk concentrate. The total processing time did not
exceed 40 minutes.

Taken together, the results demonstrate that temperature treatment
exerts a significant influence on the physicochemical characteristics of
buttermilk and on the efficiency of membrane fractionation {[}13{]}.
From the standpoint of optimizing acidity, protein, and fat retention,
30\,°C emerged as the most advantageous processing temperature.

The following table presents the physicochemical characteristics of
natural cow's milk, goat's milk, and buttermilk.
\end{multicols}
\vspace{-0.5cm}
\tcap{Table 2 - Physic-chemical parameters of buttermilk concentrate
obtained by ultrafiltration}
\vspace{-0.3cm}
\begin{longtblr}[
  label = none,
  entry = none,
]{
  width = \linewidth,
  colspec = {
    Q[120]
    Q[73]
    *{10}{Q[0.05\textwidth]}
  },
  cells = {c},
  cells = {font = \small},
  row{1} = {font=\bfseries\small},
  row{2} = {font=\bfseries\small},
  cell{1}{1} = {r=2}{},
  cell{1}{2} = {r=2}{},
  cell{1}{3} = {c=2}{},
  cell{1}{5} = {c=2}{},
  cell{1}{7} = {c=2}{},
  cell{1}{9} = {c=2}{},
  cell{1}{11} = {c=2}{},
  cell{3}{1} = {c=12}{},
  cell{9}{1} = {c=12}{0.923\linewidth},
  hlines,
  vlines,
}
\begin{sideways}Name of indicators\end{sideways} & \begin{sideways}Buttermilk control\end{sideways} & \begin{sideways}Buttermilk t=20ºC\end{sideways} & & \begin{sideways}Buttermilk t=30ºC\end{sideways} & & \begin{sideways}Buttermilk t=40ºC\end{sideways} & & \begin{sideways}Buttermilk t=50 ºC\end{sideways} & & \begin{sideways}Buttermilk t=60 ºC\end{sideways} & \\
 & & \begin{sideways}Permeate(Whey Protein)\end{sideways} & \begin{sideways}Retentate (Clarified Whey)\end{sideways} & \begin{sideways}Permeate (Whey Protein)\end{sideways} & \begin{sideways}Retentate (Clarified Whey)\end{sideways} & \begin{sideways}Permeate (Whey Protein)\end{sideways} & \begin{sideways}Retentate (Clarified Whey)\end{sideways} & \begin{sideways}Permeate (Whey Protein)\end{sideways} & \begin{sideways}Retentate (Clarified Whey)\end{sideways} & \begin{sideways}Permeate (Whey Protein)\end{sideways} & \begin{sideways}Retentate (Clarified Whey)\end{sideways} \\
n=3 & & & & & & & & & & & \\
Active acidity, pH & 7.14 ± 0.98ᵃ & 6.84 ± 0.55ᵃᵇ & 7.14 ± 0.55ᵃ & 6,63 ± 0,02ᵃ & 7,05 ± 0,15ᵃ & 6.58 ± 0.55ᵃᵇ & 6,68 ± 0,14ᵃᵇ & 6,73 ± 0,25ᵃᵇ & 6,68 ± 0,75ᵃᵇ & 6,56 ± 0,85ᵇ & 6,42 ± 0,95ᵇ \\
Titrated acidity, ̊T & 16 ± 1.2ᵃ & 16.3 ± 0.55ᵃ & 22.1 ± 0.5ᵇ & 25.6 ± 0.4ᶜ & 23.6 ± 0.22ᵇ & 26.6 ± 0.48ᶜ & 27 ± 0.48ᶜ & 30.2 ± 0.15ᵈ & 27.1 ± 0.48ᶜ & 20.66 ± 0.85ᵃᵇ & 19.1 ± 0.45ᵃ \\
Mass fraction of fats, \% & 0.69 ± 0.85ᵃ & 0.70 ± 0.11ᵃ & 0.06 ± 0.48ᵇ & 0.72 ± 0.44ᵃ & 0.08 ± 0.85ᵃ & 0.74 ± 1.02ᵃ & 0.09 ± 0.55ᵇ & 0.81 ± 0.87ᵃ & 0.10 ± 0.62ᶜ & 0.95 ± 0.75ᵃᵇ & 0.10 ± 0.62ᵇ \\
Mass fraction of proteins, \% & 2.89 ± 0.15ᵃ & 3.94 ± 1.55ᵇ & 1.90 ± 0.84ᶜ & 4.00 ± 0.84ᵇ & 1.82 ± 0.99ᶜ & 3.90 ± 0.15ᵇ & 1.44 ± 0.56ᶜ & 3.98 ± 0.55ᵇ & 1.79 ± 0.45ᶜ & 3.80 ± 0.45ᵇ & 1.76 ± 0.96ᶜ \\
Mass fraction of dry substances, \% & 10.16 ± 0.64ᵃ & 10.29 ± 0.95ᵃ & 5.29 ± 0.44ᵇ & 17.7 ± 0.44ᶜ & 5.10 ± 0.78ᵇ & 16.0 ± 0.35ᶜ & 4.06 ± 0.85ᵇ & 16.67 ± 0.42ᶜ & 5.02 ± 0.71ᵇ & 15.38 ± 0.16ᶜ & 4.94 ± 0.17ᵇ \\
The average values in a row marked with different letters (a, b, c) differ from each other statistically at p 0.05. Values with the same letters, as well as combined symbols (ab), do not differ statistically. & & & & & & & & & & & 
\end{longtblr}

\tcap{Table 3 - Physicochemical characteristics of natural cow's milk, goat's milk, and buttermilk}
\begin{longtblr}[
  label = none,
  entry = none,
]{
  width = \linewidth,
  colspec = {Q[313]Q[129]Q[125]Q[371]},
  cells = {c},
  cells = {font = \small},
  row{1} = {font=\bfseries},
  cell{1}{1} = {r=2}{},
  cell{2}{2} = {c=3}{0.625\linewidth},
  hlines,
  vlines,
}
Name of indicators                  & Сow’s milk   & Goat’s milk  & Buttermilk permeate after ultrafiltration \\
                                    & n=3          &              &                                           \\
Active acidity, pH                  & 6,2 ± 0,12   & 7,01 ± 0,11  & 6,63 ± 0,02                               \\
Titrated acidity, ̊T                & 16 ± 0,03    & 17 ± 0,01    & 25,6 ± 0,04                               \\
Mass fraction of fats, \%           & 3,98 ± 0,59  & 4,5 ± 0,01   & 0,72 ± 0,45                               \\
Mass fraction of proteins, \%       & 3,47 ± 0,12  & 3,1 ± 0,15   & 4,00 ± 0,85                               \\
Mass fraction of dry substances, \% & 10,56 ± 0,55 & 11,85 ± 0,78 & 17,7 ± 0,44                               
\end{longtblr}

\begin{multicols}{2}
As shown in Table 3, cow's milk is characterized by moderate acidity, a
fat content of about 3.98\%, and a protein level of 3.47\%, which are
within the standard range for natural cow's milk. Goat's milk shows a
slightly higher pH and fat content (4.5\%), consistent with its
naturally richer composition. The parameters of both types of milk meet
standard physicochemical requirements. The buttermilk permeate obtained
after ultrafiltration demonstrates lower fat and higher protein and dry
matter contents, reflecting the concentration effect of the membrane
process.

The buttermilk obtained through ultrafiltration underwent a
comprehensive organoleptic evaluation. Sensory attributes often exert a
more profound influence on consumer preferences than the chemical
composition or nutritional value of a product, ultimately shaping market
demand {[}14{]}. This underscores the critical importance of in-depth
sensory analysis at every stage of production.
\end{multicols}

\tcap{Table 4 - Evaluation of organoleptic parameters of buttermilk concentrates obtained by ultrafiltration, score}
\begin{longtblr}[
  label = none,
  entry = none,
]{
  cells = {c},
  cells = {font = \small},
  cell{1}{2} = {font=\bfseries},
  cell{1}{3} = {font=\bfseries},
  cell{2}{1} = {r=5}{},
  cell{7}{1} = {r=5}{},
  hlines,
  vlines,
}
            & Characteristics of organoleptic parameters                              & Score \\
Taste       & Milky, clean, moderately sweet taste with a hint of pasteurization      & 5     \\
            & Unexpressed milky, moderately sweet                                     & 4     \\
            & Empty                                                                   & 3     \\
            & The intrusive sweet taste                                               & 2     \\
            & Uncharacteristic for this product                                       & 1     \\
Consistency & Homogeneous, fluid, slightly viscous liquid without sediment and flakes & 5     \\
            & Homogeneous liquid                                                      & 4     \\
            & Excessively liquid or excessively viscous                               & 3     \\
            & Heterogeneous, with a visible protein deposit                           & 2     \\
            & Heterogeneous, floccose, stringy                                        & 1     
\end{longtblr}

\begin{multicols}{2}
In the case of fermented dairy products with a complex raw material
matrix, the development of their sensory profile is largely dictated by
the quality of the initial dairy base {[}15{]}. Among the various
parameters that define its consumer appeal, taste, aroma, and texture
stand out as the most pivotal. To quantify these attributes, a tailored
scoring system was devised specifically for buttermilk concentrates in
the Table 4.

The organoleptic evaluation revealed that the buttermilk concentrates
obtained via ultrafiltration, containing between 10\% and 17.7\% total
solids, exhibited a clean, pleasant flavor profile with a subtly sweet
undertone. Notably, the intensity of flavor increased proportionally
with the degree of concentration. Sensory attributes were assessed using
a five-point scale. The results of the organoleptic analysis of
buttermilk are presented in Table 5.
\end{multicols}

\tcap{Table 5 - Results of organoleptic analysis of buttermilk}
\begin{longtblr}[
  label = none,
  entry = none,
]{
  cells = {c},
  cells = {font = \small},
  row{1} = {font=\bfseries},
  cell{2}{1} = {c=8}{},
  hlines,
  vlines,
}
\begin{sideways}Indicators\end{sideways} & \begin{sideways}Buttermilk, 10.29\% TS\end{sideways} & \begin{sideways}Buttermilk, 10,16 TS\end{sideways} & \begin{sideways}Buttermilk, 10,29 TS\end{sideways} & \begin{sideways}Buttermilk, 16 TS\end{sideways} & \begin{sideways}Buttermilk, 16,67 TS\end{sideways} & \begin{sideways}Buttermilk, 15,38 TS\end{sideways} & \begin{sideways}Buttermilk, 17,7 TS\end{sideways} \\
n=7         &                        &                      &                      &                   &                      &                      &                     \\
Taste       & 4,5±0,12               & 4,2± 0,14            & 4,21±0,04            & 4,32±0,02         & 4,11±0,05            & 4,32±0,01            & 4,9±0,02            \\
Consistency & 4,12±0,02              & 4,23±0,05            & 4,20±0,05            & 4,32±0,07         & 4,74±0,04            & 4,8±0,01             & 4,9±0,01            
\end{longtblr}

\begin{multicols}{2}
Based on the data in the table, the best sensory results were obtained
for the buttermilk sample with 17.7\% total solids. This variant
received the highest scores for both taste and consistency, indicating a
richer flavor and a more stable, pleasant texture. The increased amount
of total solids likely improves the product's structure, making it more
balanced and more appealing in terms of sensory characteristics compared
to the other samples.

Based on a comprehensive series of analyses, the ultrafiltered
buttermilk concentrate with a 17.7\% solids content was selected as the
optimal base for the development of fermented dairy products enriched
with high-quality protein.

Numerous studies {[}16{]} have demonstrated that the rheological
behavior and syneresis of fermented milk gels are significantly
influenced by the dispersion characteristics of protein particles in the
starting material. In this context, investigating the transformations in
the structural viscosity of buttermilk - driven by its protein
components under the effect of ultrafiltration -- as well as the nature
of the resulting protein microstructures, presents both scientific and
practical value {[}17{]}.

As a dairy base meeting the defined requirements, the buttermilk
concentrate obtained through ultrafiltration proved to be the most
suitable candidate. Its compositional characteristics suggest it offers
a favorable environment for the growth of starter cultures, including
probiotic strains. Moreover, the incorporation of whey concentrate into
the dairy matrix is expected to substantially enhance the content of
high-quality whey proteins in the final product {[}18{]}.

To develop technology for protein-enriched fermented dairy products, a
tailored multicomponent mixture was formulated using skimmed cow's milk,
goat's milk, and buttermilk concentrate. Building on earlier studies
that explored the optimal ratio of cow's to goat's milk, it was
established that a 50:50 blend yielded the best functional and sensory
balance. However, a higher proportion of goat's milk significantly
increases production costs.

With this in mind, the current study sought to refine the formulation by
systematically reducing the goat's milk content while compensating with
increased levels of buttermilk concentrate. Four experimental
formulations of the dairy base were prepared using the following
proportions of cow's milk / goat's milk / buttermilk concentrate:

Variant 1: 50/40/10

Variant 2: 50/30/20

Variant 3: 50/20/30

Variant 4: 50/10/40

Each variant underwent detailed physicochemical analysis under
laboratory conditions using standard methodologies. The resulting data
are presented in Table 6.
\end{multicols}

\tcap{Table 6 - Physic-chemical parameters of a multicomponent mixture using buttermilk concentrate}
\begin{longtblr}[
  label = none,
  entry = none,
]{
  width = \linewidth,
  colspec = {Q[83]Q[253]Q[270]Q[202]Q[169]},
  cells = {c},
  cells = {font = \small},
  row{1} = {font=\bfseries},
  hlines,
  vlines,
}
Variant & {\textbf{Milk mixture ratio}\\\textbf{(cow/goat/buttermilk, \%)}} & Dry substance, \% by mass (n=3) & Proteins, g/100 g (n=3) & Fats, g/100 g (n=3) \\
1       & 50/40/10                                                          & 18.25±0,15                      & 6.01±0,02               & 3.65±0,56           \\
2       & 50/30/20                                                          & 24.00±0,18                      & 8.67±0,45               & 3.40±0,45           \\
3       & 50/20/30                                                          & 29.75±0,65                      & 11.33±1,20              & 3.15±0,15           \\
4       & 50/10/40                                                          & 35.50±0,95                      & 14.00±0,42              & 2.90±0,78           
\end{longtblr}

\begin{multicols}{2}
The results of the conducted study demonstrate a clear upward trend in
both total solids and protein content as the proportion of buttermilk
concentrate in the dairy base increases -- from Variant 1 (10\%
concentrate) to Variant 4 (40\%). Specifically, the total solids content
rose from 18.25\% to 35.50\%, while protein levels increased from 6.01\%
to 14.00\%. These findings underscore the beneficial role of buttermilk
concentrate in enhancing the density, nutritional richness, and
biological value of the final product.

It has been confirmed that buttermilk concentrate serves as a valuable
source of high-quality protein, making it a promising ingredient in the
formulation of functional food products.

{\bfseries Conclusion.} Thus, the incorporation of buttermilk concentrates
into a dairy blend containing cow's and goat's milk facilitates a
simultaneous increase in complete protein content and a reduction in
fat, offering a significant advantage in the development of functional
and dietetic dairy products.

Moreover, it was found that increasing the proportion of buttermilk
concentrate in the formulation leads to a marked rise in total solids
and a notable improvement in sensory characteristics -- such as taste,
aroma, and texture. Owing to its finely dispersed protein composition
and high biological value, buttermilk concentrate contributes to the
formation of an optimal gel structure in fermented dairy matrices and
creates a favorable environment for the growth of starter cultures,
including probiotic microorganisms.

\emph{{\bfseries Funding.}} This research was carried out within the
framework of the targeted scientific and technical program IRN
BR24892775, entitled «Development of Technology for the Integrated and
Advanced Processing of Agricultural Raw Materials for the Production of
High-Quality and Safe Food Products».
\end{multicols}

\begin{center}
{\bfseries Литература}
\end{center}

\begin{refs}
1. Claire C., Timlin M., Hogan S.\,A., Murphy E.\,G., O'Callaghan T.\,F.,
Brodkorb A., Hennessy D., Fitzpatrick E., O' Donavan
M., McCarthy K., Murphy J.\,P., Yin X., Brennan L. Impact of Dietary
Regime on the Metabolomic Profile of Bovine Buttermilk and Whole Milk
Powder// Metabolomics. - 2024. - Vol.20(5): \,93.
\href{https://doi.org/10.1007/s11306-024-02157-4}{DOI
10.1007/s11306-024-02157-4}.

2. Санакоева А.Г., Маргиева Ф.Т. Разработка технологии кисломолочного
напитка на основе концентрата пахты с сиропом облепихи// Достижения
науки - сельскому хозяйству. Материалы Всероссийской
научно-практической конференции, Владикавказ. - 2017. -T.II. -С.
238-241.

3. Национальный стандарт Российской Федерации. ГОСТ 53359-2009. Молоко и
продукты переработки молока. Метод определения pH. // Москва:
Стандартинформ. - 2009. -11 c. URL:
https://meganorm.ru/Data2/1/4293828/4293828079.pdf. -Дата обращения:
11.07.2025.

4. Государственный стандарт Российской Федерации. ГОСТ 3624-92. Молоко и
молочные продукты. Метод определения титруемой кислотности. //Москва:
Издательство стандартов. - 1992. -8 c. URL:
\url{https://files.stroyinf.ru/Data/100/10071.pdf}. -Дата обращения:
11.07.2025.

5. Государственный стандарт Российской Федерации. ГОСТ 3626-73. Молоко и
молочные продукты. Методы определения влаги и сухого вещества //
Межгосударственный стандарт. -Москва: Издательство стандартов. -1974.
-10 c. URL: \url{https://internet-law.ru/gosts/gost/36972}. -Дата
обращения: 11.07.2025.

6. Государственный стандарт Российской Федерации. ГОСТ Р ИСО
22935-2-2011. Молоко и молочные продукты. Органолептический анализ.
Часть 2. Рекомендуемые методы органолептической оценки //
Межгосударственный стандарт. -- Москва: Стандартинформю -2012. -19 с.
URL: https://files.stroyinf.ru/Data2/1/4293796/4293796164.pdf. -Дата
обращения: 11.07.2025.

7. Острецова Н. Г., Боброва А. В. Использование нанофильтрационных
концентратов пахты и сыворотки для кисломолочных продуктов с
повышенной массовой долей белка // Food systems. - 2021. - Т.4 (2). -
С.134-143. DOI
\href{https://doi.org/10.21323/2618-9771-2020-4-2-134-143}{10.21323/2618-9771-2020-4-2-134-143}.

8. Скокова О. И., Чеканова Ю. Ю., Мелех Т. В. Исследование влияния пахты
в составе комбинированной сливочной смеси на стойкость и стабильность
свойств сметаны при хранении // Сборник научных трудов «Актуальные
вопросы переработки мясного и молочного сырья». -2021. -Т.1 (15). -С.
136-145. DOI 10.47612/2220-8755-2020-15-136-145.

9. Зобкова З. С., Харитонов Д. В., Фурсова Т. П., Зенина Д. В., Гаврилина
А. Д., Шелагинова И. Р., Шефов Д. А. Влияние состава белковой фазы
кисломолочных продуктов на их реологические свойства // Переработка
молока. - 2014. - №10 (180). - С.50-51.

10. Ali A.H. Current knowledge of buttermilk: Composition, applications in
the food industry, nutritional and beneficial health characteristics
// International Journal of Dairy Technology. -2018. -Vol.72(2). -P.
169-182. DOI 10.1111/1471-0307.12572.

11. Krebs L., Bérubé A., Iung J., Marciniak A., Turgeon S.L., Brisson G.
Impact of Ultra-High-Pressure Homogenization of Buttermilk for the
Production of Yogurt//Foods. -2021. -Vol.10(8):1757. DOI
10.3390/foods10081757.

12. Национальный стандарт Российской Федерации. ГОСТ 34354-2017. Іркіт и
напитки на её основе//Москва: Стандартинформ. -2017. -19 c. URL:
\url{https://files.stroyinf.ru/Data2/1/4293740/4293740472.pdf}. - Дата
обращения: 11.07.2025.

13. Кручинин А. Г., Агаркова Е. Ю. Использование мембранных технологий при
концентрировании вторичного молочного сырья // Переработка молока. -
2017. - № 12 (218). -С.54--55.

14. Kumar R., Kaur M., Garsa A.K., Shrivastava B., Reddy V.P., Tyagi A.
Natural and cultured buttermilk //Fermented Milk and Dairy Products.
-2015. -P.203--225. ISBN 9781466577978.

15. Buey B., Bellés A., Abad I., Vergara C., Mesonero J.E., Sánchez L.,
Grasa L., et al. Comparative effect of bovine buttermilk, whey, and
lactoferrin on the innate immunity receptors and oxidative status of
intestinal epithelial cells // National Library of Medicine. -2021.
-Vol.99(1). -P.54-60. DOI 10.1139/bcb-2020-0121.

16. Ahmad N., Manzoor M.F., Shabbir U., Ahmed S., Ismail T., Saeed F.,
Nisa M., Anjum F.M., Hussain S. Health lipid indices and
physicochemical properties of dual fortified yogurt with extruded
flaxseed omega fatty acids and fibers for hypercholesterolemic
subjects // Food Science \& Nutrition. -2019. -Vol.8(1). -P.273-280.
DOI 10.1002/fsn3.1302.

17. Barukčić I., Lisak Jakopović K., Božanić R. Valorisation of whey and
buttermilk for production of functional beverages: an overview of
current possibilities // Food Technology and Biotechnology. -2019.
-Vol.57(4). -P.448-460. DOI
\href{https://doi.org/10.17113/ftb.57.04.19.6460}{10.17113/ftb.57.04.19.6460}.

18. Абделлатыф С.С., Тихомирова Н. А. Пахта: один из источников молочных
минорных компонентов // Пищевые ингредиенты России 2019: материалы
конференции. -2019. -С.6-9.
\end{refs}

\begin{center}
{\bfseries References}
\end{center}

\begin{refs}
1. Claire C., Timlin M., Hogan S.\,A., Murphy E.\,G., O'Callaghan
T.\,F., Brodkorb A., Hennessy D., Fitzpatrick E.,
O' Donavan M., McCarthy K., Murphy J.\,P., Yin X.,
Brennan L. Impact of Dietary Regime on the Metabolomic Profile of Bovine
Buttermilk and Whole Milk Powder// Metabolomics. -2024. -Vol.20(5):
\,93. \href{https://doi.org/10.1007/s11306-024-02157-4}{DOI
10.1007/s11306-024-02157-4}.

2. Sanakoeva A.G., Margieva F.T. Razrabotka tekhnologii kislomolochnogo
napitka na osnove kontsentrata pakhty s siropom oblepikhi //
Dostizheniya nauki - sel' skomu khozyaystvu. Materialy
Vserossiyskoy nauchno-prakticheskoy konferentsii, Vladikavkaz. -2017.
-T. II. -S.238-241. {[}in Russiаn{]}

3. Nacional' nyj standart Rossijskoj Federacii. GOST
53359-2009. Moloko i produkty pererabotki moloka. Metod opredelenija
pH.//Moskva: Standartinform. -2009. -11 s. URL:
https://rosgosts.ru/file/gost/67/100/gost\_r\_53359-2009.pdf. -Date of
access: 11.07.2025. {[}in Russiаn{]}

4. Gosudarstvennyj standart Rossijskoj Federacii. GOST 3624-92. Moloko i
molochnye produkty. Metod opredelenija titruemoj kislotnosti. //Moskva:
Izdatel' stvo standartov. -1992. -8 s. URL:
https://files.stroyinf.ru/Data/100/10071.pdf. -Date of access:
11.07.2025. {[}in Russiаn{]}

5. Gosudarstvennyj standart Rossijskoj Federacii. GOST 3626-73. Moloko i
molochnye produkty. Metody opredelenija vlagi i suhogo veshhestva //
Mezhgosudarstvennyj standart. -Moskva: Izdatel' stvo
standartov. - 1974. - 10 s. URL:
https://internet-law.ru/gosts/gost/36972. -Date of access: 11.07.2025.
{[}in Russiаn{]}

6. Gosudarstvennyy standart Rossiyskoy Federatsii. GOST R ISO
22935-2-2011. Moloko i molochnye produkty. Organolepticheskiy analiz.
Chast' 2. Rekomenduemye metody organolepticheskoy otsenki //
Mezhgosudarstvennyy standart. -- Moskva: Standartinform. -2012. -19 s.
URL: \url{https://files.stroyinf.ru/Data2/1/4293796/4293796164.pdf}.
-Date of access: 11.07.2025. {[}in Russiаn{]}

7. Ostrecova N. G., Bobrova A. V. Ispol' zovanie
nanofil' tracionnyh koncentratov pahty i syvorotki dlja
kislomolochnyh produktov s povyshennoj massovoj dolej belka // Food
systems. - 2021. - T.4 (2). - S.134--143. DOI
10.21323/2618-9771-2020-4-2-134-143. {[}in Russiаn{]}

8. Skokova O. I., Chekanova Ju. Ju., Trilinskaja E. A., Avtushenko V.
V., Meleh T. V. Issledovanie vlijanija pahty v sostave kombinirovannoj
slivochnoj smesi na stojkost'{} i
stabil' nost'{} svojstv smetany pri
hranenii//Sbornik nauchnyh trudov «Aktual' nye voprosy
pererabotki mjasnogo i molochnogo syr' ja». -2021. -T.
15. - S.136-145. DOI 10.47612/2220-8755-2020-15-136-145. {[}in
Russiаn{]}

9. Zobkova Z. S., Haritonov D. V., Fursova T. P., Zenina D. V.,
Gavrilina A. D., Shelaginova I. R., Shefov D. A. Vlijanie sostava
belkovoj fazy kislomolochnyh produktov na ih reologicheskie svojstva //
Pererabotka moloka. - 2014. - № 10. - S.50--51.

10. Ali A.H. Current knowledge of buttermilk: Composition, applications
in the food industry, nutritional and beneficial health characteristics
// International Journal of Dairy Technology. -2019. -Vol.72(2). -P.
169-182. DOI 10.1111/1471-0307.12572.

11. Krebs L., Bérubé A., Iung J., Marciniak A., Turgeon S.L., Brisson G.
Impact of Ultra-High-Pressure Homogenization of Buttermilk for the
Production of Yogurt//Foods. -2021. -Vol.10(8):1757. DOI
10.3390/foods10081757.

12. Nacional' nyj standart Rossijskoj Federacii. GOST
34354-2017. Pahta i napitki na ejo osnove. // Moskva: Standartinform.
-2017. -19 c. URL:
https://files.stroyinf.ru/Data2/1/4293740/4293740472.pdf. -Date of
access: 11.07.2025. {[}in Russiаn{]}

13. Kruchinin A. G., Agarkova E. Ju. Ispol' zovanie
membrannyh tehnologij pri koncentrirovanii vtorichnogo molochnogo
syr' ja // Pererabotka moloka. - 2017. - № 12. - S.
54-55. {[}in Russiаn{]}

14. Kumar R., Kaur M., Garsa A.K., Shrivastava B., Reddy V.P., Tyagi A.
Natural and cultured buttermilk //Fermented Milk and Dairy Products.
-2015. -P.203--225. ISBN 9781466577978.

15. Buey B., Bellés A., Abad I., Vergara C., Mesonero J.E., Sánchez L.,
Grasa L., et al. Comparative effect of bovine buttermilk, whey, and
lactoferrin on the innate immunity receptors and oxidative status of
intestinal epithelial cells // National Library of Medicine. -2021.
-Vol.99(1). -P.54--60. DOI 10.1139/bcb-2020-0121.

16. Ahmad N., Manzoor M.F., Shabbir U., Ahmed S., Ismail T., Saeed F.,
Nisa M., Anjum F.M., Hussain S. Health lipid indices and physicochemical
properties of dual fortified yogurt with extruded flaxseed omega fatty
acids and fibers for hypercholesterolemic subjects // Food Science \&
Nutrition. -2020. -Vol.8(1). -P.273-280. DOI 10.1002/fsn3.1302.

17. Barukčić I., Lisak Jakopović K., Božanić R. Valorisation of whey and
buttermilk for production of functional beverages: an overview of
current possibilities // Food Technology and Biotechnology. -2019. -Vol.
57(4). -P.448-460. DOI 10.17113/ftb.57.04.19.6396.

18. Abdellatyf S.S., Tihomirova N. A. Pahta: odin iz istochnikov
molochnyh minornyh komponentov // Pishhevye ingredienty Rossii 2019:
materialy konferencii. -2019. -S.6-9. {[}in Russiаn{]}
\end{refs}

\begin{info}
\hspace{1em}\emph{{\bfseries Information about authors}}

Tultabaeva Т.Ch. - Doctor of Technical Sciences, Professor, S.Seifullin
Kazakh Agro Technical Research University, Astana, Kazakhstan, e-mail:
tamara\_tch@list.ru;

Zhakupova G.N. - Candidate of Technical Sciences, Professor, S.Seifullin
Kazakh Agro Technical Research University, Astana, Kazakhstan, e-mail:
gulmira-zhak@mail.ru;

Makangali К.К.- PhD, Аssociate Professor, S.Seifullin Kazakh Agro
Technical Research University, Astana, Kazakhstan, e-mail:
k.makangali@mail.ru;

Sagandyk A.T. - PhD, S.Seifullin Kazakh Agro Technical Research
University, Astana, Kazakhstan, e-mail:
assema.bukeyeva@gmail.com;

Muldasheva A.H. - Master of Technical Sciences, S.Seifullin Kazakh Agro
Technical Research University, Astana, Kazakhstan, e-mail:
aknurmuldasheva@gmail.com;

Dairova K.Sh - Master's student, S.Seifullin Kazakh Agro Technical
Research University, Astana, Kazakhstan, e-mail:
dairova111@mail.ru.

\hspace{1em}\emph{{\bfseries Сведения об авторах}}

Тултабаева Т.Ч. - д.т.н. профессор, Казахский агротехнический
исследовательский университет имени С. Сейфуллина, Астана, Казахстан,
e-mail: tamara\_tch@list.ru;

Жакупова Г.Н. - к.т.н., профессор, Казахский агротехнический
исследовательский университет имени С. Сейфуллина, Астана, Казахстан,
e-mail:
gulmira-zhak@mail.ru;

Макангали К.К.- доктор PhD, ассоциированный профессор, Казахский
агротехнический исследовательский университет имени С. Сейфуллина,
Астана, Казахстан, e-mail:
k.makangali@mail.ru;

Сагандык А.Т. - доктор PhD, Казахский агротехнический исследовательский
университет имени С. Сейфуллина, Астана, Казахстан, e-mail:
assema.bukeyeva@gmail.com;

Мулдашева А.Х. - магистр технических наук, Казахский агротехнический
исследовательский университет имени С. Сейфуллина, Астана, Казахстан,
e-mail: aknurmuldasheva@gmail.com;

Даирова К.Ш. - магистрант, Казахский агротехнический исследовательский
университет имени С. Сейфуллина, Астана, Казахстан, e-mail:
dairova111@mail.ru.
\end{info}
